(1) 将 Schrödinger 波函数展开为球谐函数，径向约化函数 \(u_l(r)=rR_l(r)\) 在球外满足 \(u_l''+[k^2-l(l+1)/r^2]u_l=0\)，在硬球表面需要 \(u_l(a)=0\)。低能 \(ka\ll1\) 时高 \(l\) 受离心势垒抑制，主导为 \(l=0\)。该分波外部方程是 \(u_0''+k^2u_0=0\)，远区可写 \(u_0\propto\sin(kr+\delta_0)\)。

(2) 施加 \(u_0(a)=0\)，连续选择低能支得到 \(\delta_0=-ka\)（模 \(\pi\)）。部分波振幅公式 \(f_0(k)=(e^{2i\delta_0}-1)/(2ik)=e^{i\delta_0}\sin\delta_0/k\) 给 \(f_0(k)=-e^{-ika}\sin(ka)/k\)，在 s 波近似下与散射方向无关。写作 \(-a\) 只是低能极限，有限 \(ka\) 仍有振荡相位与 \(\sin(ka)/(ka)\) 修正。

(3) 对角向常数的 s 波振幅，微分截面 \(d\sigma/d\Omega=|f_0|^2=\sin^2(ka)/k^2\)，积分得 \(\sigma_{0}=4\pi\sin^2(ka)/k^2\)。当 \(ka\to0\) 时 \(\sin(ka)\sim ka\)，故 \(\sigma_{\rm tot}=4\pi a^2[1-(ka)^2/3+\cdots]\)，其中这里的总截面表达只保留 s 波；高分波在更高阶会影响完整修正系数。因此严格的主导项是 \(4\pi a^2\)，不能把所写次阶项无条件宣称为完整硬球散射截面的精确展开。

(4) 零能 \(l=0\) 径向方程为 \(u_0''=0\)，外部解由硬球边界给 \(u_0(r)\propto r-a\)。将其与零能散射长度定义 \(u_0(r)\propto r-a_s\) 对比，得 \(a_s=a>0\)。低能振幅 \(f_0\to-a_s\)，总截面 \(4\pi a_s^2=4\pi a^2\)。经典几何碰撞截面为 \(\pi a^2\)，量子低能波长远大于硬球时因 s 波相干散射呈四倍系数；这不是违反几何守恒，而是波散射定义与经典直线轨道截面的差异。
